\section{Evaluating Structural Discrimination}
\label{sec:experimental-protocol}

No single accuracy score can answer this question. A representation may
distinguish two graphs even when a chosen classifier cannot exploit the
difference; conversely, a classifier may succeed by using graph size rather
than the intended structure. We therefore separate three questions: does the
representation change, can a simple model use the change, and what does the
change cost?

\subsection{Node and Graph Feature Construction}

An operator expression still has to become a numerical representation before a
predictive model can use it. For expression $P$, identity width $b$, and
diagnostic model $f$, this proceeds as
\begin{equation}
 G\xrightarrow{P}A_P(G)\xrightarrow{v_b}x_{P,b}(G)
 \xrightarrow{f}\widehat{y}.
\end{equation}
The native representation is node-level. For base node $v$ and feature
bucket $j$, let
\begin{equation}
 Z_{P,b}(G)_{v,j}=\sum_{u\in V_I}
 \mathbf{1}[v\in V_{\mu(u)}]
 \mathbf{1}[\phi_b(\mu(u))=j]a(u).
\end{equation}
A row records the interpreted structures in which a node occurs. Graph-level
features are obtained by sum pooling,
\begin{equation}
 x_{P,b}(G)=\sum_{v\in V}Z_{P,b}(G)_{v,:}.
\end{equation}
This creates an exact node-to-graph consistency relation, but it is not always
a raw occurrence histogram. With unit contributions,
\begin{equation}
 x_{P,b,j}(G)=\sum_{u\in V_I}|V_{\mu(u)}|
 \mathbf{1}[\phi_b(\mu(u))=j],
\end{equation}
so pooling counts node--structure incidences and weights an occurrence by its
mapped-subgraph size. It equals the raw occurrence count if mapped components
are singletons or if $a(u)=1/|V_{\mu(u)}|$. We therefore compare raw
occurrences with incidence pooling.

Columns zero and one are reserved:
\begin{equation}
 Z_{v,0}=1,\qquad Z_{v,1}=\deg_G(v).
\end{equation}
After pooling they equal $|V|$ and $2|E|$, respectively, for undirected
graphs and for directed total degree (in-degree plus out-degree). These
features can create size shortcuts. All controlled results are therefore
reported with and without them, and topology-focused pairs are matched on
$|V|$ and $|E|$.

\subsection{Intrinsic Discrimination and Predictive Accessibility}

We first ask whether two graphs receive different signatures, without fitting
a model. For verified non-equivalent pairs \(\mathcal{Q}\), intrinsic
discrimination is
\begin{equation}
 D_{\mathrm{int}}(P,b)=\frac{1}{|\mathcal{Q}|}
 \sum_{(G,H)\in\mathcal{Q}}
 \mathbf{1}[\Sigma_{P,b}(G)\ne\Sigma_{P,b}(H)].
\end{equation}
Results are stratified by the changed factor---including cycles, paths, rays,
attachment patterns, and labels---rather than reduced to one score.

Different signatures do not guarantee that a task-relevant distinction is
easy to recover. Simple supervised models test accessibility rather than state-of-the-art
prediction: logistic or multinomial logistic regression for categories, ridge
regression for counts, and a linear SVM as a margin-based sensitivity analysis.
If signatures and linear probes disagree, we report the distinction without
attributing it to representation collapse. Feature vocabularies are fitted on
training data only.

\subsection{Controlled Structural Contrasts and Baselines}

To know why a representation succeeds, we need graph pairs whose relevant
difference is known. The primary synthetic family composes cycles, paths, and star-like rays while
recording its generating parameters. Factors include cycle count and length,
path length, ray count and length, attachment arrangement, expression depth,
and node- and edge-label alphabets. Each graph is paired with identifier
permutations that must remain equivalent and with single-factor interventions
that must first be verified as non-isomorphic under the declared attributes.
Hard pairs additionally match node count, edge count, degree distribution,
density, and label histogram. Operator expressions and pair-selection rules
are frozen before confirmatory evaluation. Generator parameter ranges and
structural combinations held out during design test whether conclusions extend
beyond the examples used to choose the expressions. ZINC supplies natural
feature frequencies and scaling measurements, not causal structural contrasts.

Controlled examples are useful only if the comparison methods are credible.
Node, edge, path, cycle, graphlet, WL, and NSPDK representations are evaluated
as expressions in the same algebra. For each method we separate decomposition,
component identity, and aggregation. Claimed baseline equivalence is checked
against an independent implementation using feature multiplicities,
graph-pair decisions, and kernel or vector values where applicable. For every
composite $P_2\circ P_1$, we compare four representations: each atomic
component, lossless concatenation of their separately vectorized signatures,
and graph-valued composition before vectorization. Concatenation is the
critical control: a composite is scientifically informative only when its
discrimination, localization, or cost differs from retaining the final atomic
features side by side.

For each separated pair we also retain a mapped witness: the interpretation
nodes whose identities differ and the corresponding base nodes and edges.
Witness correctness is checked against the known synthetic intervention. This
tests whether the mapping supports structural localization rather than merely
accompanying a graph-level signature.

\subsection{Separating Structural Collapse from Hash Collisions}

Two distinct structures can disappear for different reasons. An operator may
never expose their difference, or a later identity map may send both to the
same bucket. With $b=\texttt{nbits}$, structural hashes occupy $B=2^b-2$ buckets because
the first two columns are reserved. We distinguish four phenomena: operator
collapse, certificate collapse, bounded-hash collision, and a collision that
is statistically irrelevant to a task. Against a collision-free certificate
registry, we report distinct certificates, occupied buckets, colliding pairs,
frequency-weighted contamination, predictive changes, memory, and provenance
ambiguity over several values of $b$.

If certificate frequencies $p_i$ are heavy-tailed, the collision mass
\begin{equation}
 C_{\mathrm{mass}}(b)=\sum_{i<j}p_i p_j
 \mathbf{1}[h_b(c_i)=h_b(c_j)]
\end{equation}
may remain small even when raw collision counts are large. For a frequent
identity $i$, expected relative contamination is approximately
$ (1-p_i)/(Bp_i) $. This motivates the hypothesis that many collisions pair
rare tail identities with more common ones and have limited aggregate effect;
it does not assume that rare features are unimportant or that interpretability
survives the collision.

\subsection{Discrimination--Complexity Trade-offs}

An expression that distinguishes more pairs is not automatically preferable if
the gain requires prohibitive time or memory. Every configuration records wall time, peak memory, number of interpretation
nodes and edges, sparse nonzeros, dimensionality, serialized size, and failure
boundaries. We report separate factor-specific discrimination--runtime,
discrimination--memory, and discrimination--size Pareto frontiers. Aggregate
discrimination is a secondary summary because it can hide incomparable
profiles. Raw costs
remain primary. Dominated and negative configurations are retained because
closure permits an expression to be valid even when its added computation
provides no useful discrimination.
